\chapter{项目背景与总体方案}

\section{问题背景}

一门典型的操作系统实验课通常围绕 xv6 等教学内核逐步展开。学生先阅读教材和已有代码，再完成系统调用、进程调度、虚拟内存、文件系统与设备驱动等实验。每个 Lab 看似只要求实现一个功能，实际还需要理解相关机制，找到正确的修改位置，完成构建和 QEMU 运行，并根据测试失败继续调试。教师既关心程序是否通过测试，也关心学生能否解释自己的实现选择、关键数据结构和异常情况。

生成式 AI 能够阅读代码、解释错误并生成补丁，为学生提供更及时的帮助。但如果直接把通用编码 Agent 接入实验仓库，AI 很容易把完成代码当成唯一目标。它可能在学生尚未理解机制时给出整段实现，修改当前 Lab 以外的文件，依据不完整日志宣布测试通过，或者留下学生无法在答辩中解释的代码。最终结果也许能够运行，却没有形成从理解、设计、实现到验证的学习过程。具体表现为以下四类问题。

\begin{enumerate}
  \item \textbf{理解被答案替代：}AI 直接给出实现后，学生无法解释代码采用的机制和设计选择。
  \item \textbf{实验范围失控：}为了修好当前功能，AI 可能同时改动后续 Lab 的代码、测试程序或构建脚本，使结果难以判断。
  \item \textbf{运行结果难以复查：}模型可以声称``测试通过''，但代码仍可能存在问题，日志中也可能留有尚未处理的错误。
  \item \textbf{学生责任变得模糊：}如果设计、编码和验收都由模型完成，课程便无法区分学生自己的理解与 AI 在实现上提供的帮助。
\end{enumerate}

VeriSpecOSLab 的目标是把 AI 变成受课程流程约束的实验助手。学生先讨论并亲手写明系统设计和当前任务，再让 AI 在明确的代码范围内参与实现。平台记录实际修改、构建过程和测试结果，规格由学生提交，真实硬件结果由教学团队复核。AI 因而能够参与实验，同时保留学生应完成的理解、选择和解释。

\section{总体方案}

VeriSpecOSLab 把一次操作系统实验拆成六个前后相接的步骤。前一步必须留下明确结果，后一步才能继续，学生不会从一句自然语言要求直接跳到 AI 生成的代码。

\begin{center}
\textbf{明确任务} $\rightarrow$ \textbf{写清设计} $\rightarrow$ \textbf{限定实现} $\rightarrow$ \textbf{构建运行} $\rightarrow$ \textbf{验证记录} $\rightarrow$ \textbf{报告提交}
\end{center}

\begin{enumerate}[leftmargin=2.2em,itemsep=0.3em]
  \item \textbf{明确任务。} 学生阅读当前 Lab，可以通过 \code{agent ask} 讨论体系结构、内核组织、保护机制、模块划分和目标硬件，再说明准备解决的问题。
  \item \textbf{写清设计。} 学生亲手记录系统方案、模块职责、接口、目标和测试方法。\code{spec lint} 检查结构与引用，\code{agent review} 给出只读评审。学生修改后手动提交。
  \item \textbf{限定范围实现。} 学生先提交当前工作，平台再创建独立的临时工作目录。实现助手只能修改当前任务声明的代码范围。平台会检查 Git 实际记录的变化，拒绝无关修改。
  \item \textbf{构建并运行。} 实现助手同时交付代码和配套测试。平台校验测试描述，在临时目录中执行构建、公开测试、契约测试、固定种子 fuzz 和有限 trace 测试。Host、QEMU 和开发板运行使用项目中预先声明的步骤。
  \item \textbf{保存验证记录。} 每次执行都记录代码版本、实际命令、完整日志、输出文件和判定结果。QEMU 与真实开发板结果分开保存，开发板结果仍由人工确认。
  \item \textbf{报告与提交。} 报告直接汇总提交记录、规格和测试结果。提交包绑定同一版代码与验证记录，不依赖 AI 再写一遍``已经完成''的说明。
\end{enumerate}

课程资料和参考代码可以导入本地知识库。问答助手需要给出资料来源，学生能够回到教材或代码核对答案。

这一方案与 SYSSPEC 等规格驱动生成系统都强调先写清规格再生成代码\cite{sysspec}。VeriSpecOSLab 将这一思路用于课程实验，要求学生在每个 Lab 中完整保留设计决定、代码修改、测试过程和人工确认。

\section{核心贡献}

本项目的核心贡献是实现 VeriSpecOSLab 平台，并用一套完整源码的 RISC-V xv6 课程案例验证它能够承载真实操作系统实验。平台将项目初始化、概念问答、学生手写规格、只读评审、AI 协作实现、多类测试、过程记录、报告生成和提交归档连接成一条可执行流程。xv6 案例提供连续的 Lab 历史、模块与接口规格、QEMU 配置和 VisionFive 2 移植内容，同一份课程工程能够走通平台的主要环节。workspace 类型检查、构建、分模块自动测试和课程历史审计为平台实现提供了验证。
